% created by Uwe Schadewald
% modified by Mathias Kuntze and Ahmet Uysal
% Add, handout to documentclass arguments for condensed pdf
\documentclass[presentation, 8pt, mathserif, t]{beamer} % , aspectratio=169
\usepackage[english]{babel}
\usepackage{pgf,graphicx}
\usepackage{amsmath, amssymb}
\usepackage[utf8]{inputenc}
\usepackage{lmodern}
\usepackage{palatino}
\usepackage{multimedia}
\usepackage{pgfpages} 
\usepackage{tikz}
\usepackage{datetime}
\pdfoptionpdfminorversion=5

\usepackage{caption}
\usepackage{subcaption}
% if else
\usepackage{ifthen}
% extend table options
\usepackage{tabularx} 
\usepackage{booktabs}
\usepackage{multicol}
\usepackage{multirow}
\usepackage{eso-pic}  % package to set background image
\usepackage[calc]{picture}

% Packages and stuff for ToDo list like itempoints
\usepackage{pifont}
\newcommand{\cmark}{\ding{51}}%
\newcommand{\xmark}{\ding{55}}%
\newcommand{\open}{$\square$}
\newcommand{\done}{\rlap{$\square$}{\raisebox{1pt}{\large\hspace{1.5pt}\cmark}}\hspace{-2.5pt}}
\newcommand{\wontfix}{\rlap{$\square$}{\raisebox{1pt}{\large\hspace{1.5pt}\xmark}}}
\newcommand{\notsure}{\rlap{$\square$}{\raisebox{0.8pt}{\large\hspace{1.5pt}\textbf{?}}}}



% side bar and footer
\setbeamertemplate{headline}{	
	\leavevmode
	\vspace{-4em}	
	\hbox{		
		\begin{beamercolorbox}[wd=0.85\paperwidth,ht=10ex,dp=8ex,center]{}%			
			% navigation with subsections as dots
			\hspace{3.5em}\insertnavigation{0.7\paperwidth}{\hskip0pt plus1fill} % add navigation in footer						
			% navigation with sections, no subsections
			% \insertsectionnavigationhorizontal{0.6\paperwidth}{\hskip0pt plus1fill}{} \\ % add navigation in footer}
			
		\end{beamercolorbox} 				
	}
	\vskip0pt
}


\setbeamertemplate{footline}{	
	\leavevmode
	\vspace{-3em}
	\hbox{
		\begin{beamercolorbox}[wd=.33\paperwidth,ht=2.25ex,dp=1ex,left]{author in head/foot}%
			\hspace{5em}
			\insertshortauthor
		\end{beamercolorbox}
		\begin{beamercolorbox}[wd=.33\paperwidth,ht=2.25ex,dp=1ex,center]{title in head/foot}%
			\insertshorttitle \ - \insertshortsubtitle
		\end{beamercolorbox}	
		\begin{beamercolorbox}[wd=0.30\paperwidth,ht=10ex,dp=8ex,right]{pagenumber in head/foot}			 	
			\insertframenumber % add page numbers
		\end{beamercolorbox}
	}			
	\vskip0pt
}



\setbeamertemplate{frametitle}{
	\ifthenelse{\equal{\insertframesubtitle}{}}{
		\vspace{0.6cm}
		\huge{\insertframetitle}
	}{
		\vspace{0.6cm}
		\small{\insertframetitle}\\
		\vspace{0.3cm}
		\huge{\insertframesubtitle}
    }		
}

	
% enumerate sections
\setbeamertemplate{section in head/foot}{\hfill\insertsectionheadnumber.~\insertsectionhead}
%\setbeamertemplate{section in head/foot shaded}{\color{structure!50}\hfill\insertsectionheadnumber.~\insertsectionhead}
\setbeamertemplate{section in toc}{\inserttocsectionnumber.~\inserttocsection}

%enumerate subsections
\setbeamertemplate{subsection in head/foot}{\hfill\insertsubsectionheadnumber.~\insertsubsectionhead}
\setbeamertemplate{subsection in head/foot shaded}{\color{structure!50}\hfill\insertsubsectionheadnumber.~\insertsubsectionhead}
%\setbeamertemplate{subsection in toc}[subsections numbered]
\setbeamertemplate{subsection in toc}{\vskip0.5em\leftskip=2em\inserttocsubsection\par}

%--------------------------Common------------------------------------------------------
\setbeamercovered{transparent} % make the beamer theme invisible
\usefonttheme{structurebold}
\beamertemplatenavigationsymbolsempty % set navigations helper function to off
\setbeamertemplate{bibliography item}[text]
\setbeamertemplate{note page}[plain]

%\setlist[itemize,1]{label={$\bullet$}} % \item are using bullets
\setbeamertemplate{itemize items}[circle]
	
	

	
% create a new command to show it on two screens
% I'm using dspdfviewer.
\newcommand{\setDualView} {
	\setbeameroption{show notes on second screen=right}
}

%\AtBeginSection[]{\subsection{}}
\newcommand{\addcite}[1]{%
	\AddToShipoutPictureFG*{%
		\AtPageLowerLeft{%
			\put(0.90\paperwidth,5em){											
				\tiny{
					\cite{#1} 
				}			
			}
		}
	}	
}

% insert a frame with references -> use bibtex
\newcommand{\insertReferenceFrame}[3]{%
	\section{#1}
	\begin{frame}[allowframebreaks]
		\frametitle{#1}
		\bibliographystyle{#2}
		\bibliography{#3}
	\end{frame}	
}

\AtBeginSection[]{\subsection{}}
	





\usepackage{../KU-Beamer-Template/style/koc} 
\usepackage{minted}
\usepackage{upquote}
 
\title{KOLT Python} 
\subtitle{Review 0} 
\newdate{date}{20}{02}{2019}
\date{\displaydate{date}}
\author{Ahmet Uysal}

\titlegraphic{\includegraphics[scale=0.2]{../KU-Beamer-Template/style/images/logo_kolt.eps}}

\setbeamercovered{invisible} % transparent

\begin{document}
  \maketitle

	\frame{\frametitle{Agenda}\tableofcontents}

  \section{Recap}
  
    \begin{frame}{Installations}
      \LARGE
      \begin{itemize}
        \item Any problems at installing Python?
        \pause
        \item Are you able to launch Python from terminal/command line?
        \pause
        \item VS Code/Any editor?
        \pause
        \item Did you create and run a Python file using VS Code?
      \end{itemize}
    \end{frame}

    \begin{frame}{Comments}
			\LARGE
			\inputminted[frame=single,framesep=2pt]{python3}{../Lecture1/code-examples/comments.py}
			Python will basically ignore comments, they are purely written \textbf{for humans}!
    \end{frame}
    
    \begin{frame}{Variables}
			\LARGE
			\begin{table}[]
				\resizebox{0.8\textwidth}{!}{
				\begin{tabular}{|l|l|l|}
				\hline
				\textcolor{koc}{\textbf{Type}} &\textcolor{koc}{\textbf{Explanation}} & \textcolor{koc}{\textbf{Examples}} \\ \hline
				\textbf{\texttt{int}}  & represent \textbf{integers} & 3, 4, 17, -10 \\ \hline 
				\textbf{\texttt{float}} & represent \textbf{real numbers} & 3.0, 1.11, -109.123123 \\ \hline
				\textbf{\texttt{bool}} & represent \textbf{boolean} truth values & \texttt{True}, \texttt{False} \\ \hline
				\textbf{\texttt{str}} & A sequence of characters. & \textquotesingle Hello\textquotesingle, \textquotesingle \textquotesingle, \textquotesingle 3\textquotesingle \\ \hline
				\textbf{\texttt{NoneType}}& special and has one value, None & \texttt{None} \\ \hline
				\end{tabular}}
			\end{table}

      \begin{itemize}
        \item How to create a variable?        
        \pause
        \texttt{variable\_name = value}
        \pause
        \item How about types?
        \pause
         use \texttt{type()}
        \pause
        \item Can a variable change type?
        \pause
        \textbf{Yes!} Just assing a new value with any type.
        \pause
        \item What if we if want to convert a value between types, i.e, \textquotesingle 2\textquotesingle $\to$ 2
      \end{itemize}
    \end{frame}

    \begin{frame}{Casting}
      \LARGE
      \begin{itemize}
        \item \texttt{int(\textquotesingle 2\textquotesingle)} $\to$ 2
        \pause
        \item Any possible reasons for casting?
        \pause
         -taking user input
        \pause
         -reading numbers from a file?
        \pause
        \item Can we cast every value to every type?
        \pause
         \textbf{NO!} try \texttt{int(\textquotesingle hello\textquotesingle)}
      \end{itemize}
    \end{frame}

    \begin{frame}{Console I/O(Input/Output)}
			\huge
			\textbf{\texttt{print(*args, sep=\textquotesingle \ \textquotesingle, end=\textquotesingle \textbackslash n\textquotesingle )}}
			\begin{itemize}
				\LARGE
				\item Can take arbitrary number of arguments
				\item Separates elements with space by default
				\item Adds newline character \texttt{\textquotesingle \textbackslash n\textquotesingle} to end by default
			\end{itemize}
			
			\textbf{\texttt{input([prompt])}}
			\begin{itemize}
				\LARGE
				\item Prints the prompt to Console
				\item Program is paused until user enters something
				\item \textbf{returns an \texttt{str} object!} 
			\end{itemize}
    \end{frame}
    
  \section{Assignment Submissions}
    
    \begin{frame}{OK}
      \begin{center}
        \includegraphics[width=0.2\textwidth]{images/okpy.png}
      \end{center}
      \pause
      \begin{itemize}
        \item We will use a website called \href{https://okpy.org/}{\underline{\textit{OK(okpy.org)}}}.
        \pause
        \item Developed by UC Berkeley teachers, using Python :)
        \pause
        \item We already registered you, go to the website and login using your KU account.
        \pause
        \item You will submit your assignments using this website, you will be able work with pairs in some assignments.
        \pause
        \item Using this website we will give you personal feedbacks based on your submissions.
        \pause
        \item Let's see a demo!
      \end{itemize}
    \end{frame}

  \section{HackerRank Contests}

    \begin{frame}{Online Contests}
      \begin{itemize}
        \item Did you created an account?
        \pause
        \item How many of you looked at problems?
        \pause
        \item Any problems/questions?
      \end{itemize}
    \end{frame}

\end{document}